\documentclass{article}
\usepackage{amsmath,amssymb}
\begin{document}
Research notes for the matching-complex question.

Counterexample in rank \(3\):
\[
w=\bar a a\, b\bar b\, \bar a a\, c\bar c\, \bar a a.
\]
It is freely trivial by adjacent cancellations.  It is quasi-reduced
because the only adjacent inverse pairs are
\(\bar a a\), \(b\bar b\), \(\bar a a\), \(c\bar c\), \(\bar a a\);
after each such pair the next letter, when present, is not the first
letter of the pair.  The \(b\bar b\) and \(c\bar c\) pairs serve as
separators which prevent the shorter alternating word
\(\bar a a\bar a a\bar a a\) from violating quasi-reducedness.

Interpretation of the matching condition used throughout: compactify the
upper half-plane to a disk, put vertices at crossings, and regard
\(M\) as a finite planar graph \(\Gamma(M)\) whose boundary vertices are
the marked endpoints and whose interior vertices are transverse
crossings.  The condition in the problem says that each component of
\(M\cap\partial\overline R\) for a complementary region \(R\) is an arc
with both endpoints among the marked boundary points, and those labels
are inverse.

Structural lemmas.

1.  \(\Gamma(M)\) is a forest for every admissible matching.  Any graph
cycle is interior since boundary vertices are leaves.  Choose an
innermost cycle.  If there were graph inside its disk, then a component
inside either contains a smaller cycle or is a tree.  A wholly interior
tree would have an illegal interior leaf; a tree attached to the chosen
cycle in two places gives a smaller cycle together with a subarc of the
chosen cycle; and a tree attached in one place again has an interior
leaf.  Hence the cycle itself is a boundary component of a complementary
region, giving a component of \(M\cap\partial\overline R\) with no
marked endpoints, contrary to the definition.

2.  If a letter \(x\) and its inverse each occur once, then their
endpoints are joined by an uncrossed interval.  Let the endpoints be
\(p,q\).  The two local sides of the incident edge at \(p\) must both
belong to face-boundary arcs ending at \(q\).  In the tree component of
\(\Gamma(M)\), take the unique path from \(p\) to \(q\).  At the first
crossing on that path, the two local sides of the incoming edge leave
along two different half-edges; after deleting the crossing vertex these
half-edges lie in different tree components, so at most one can contain
\(q\).  Contradiction.  Hence no crossing lies on the path, which is a
single interval.

3.  Applying (2) to \(b,\bar b\) and \(c,\bar c\) forces the intervals
\((3,4)\) and \((7,8)\).  Since the endpoints are adjacent, each
interval is boundary-parallel and cuts off an empty disk.  Deleting
these forced intervals identifies \(F_w\) cell-by-cell with the
analogous labelled-boundary complex for
\[
u=\bar a a\bar a a\bar a a
\]
on six endpoints.  The auxiliary word \(u\) need not be quasi-reduced;
it is just the residual labelled boundary problem.

Six-endpoint computation for \(u\).

Odd endpoints are labelled \(\bar a\), even endpoints \(a\).  There are
15 pairings.  In an admissible matching, two intervals cannot cross
twice (bigon cycle), and all three pairs cannot cross (triangle cycle).
Together with mod-2 crossing parity for chords in a disk, this means
the crossing pattern is determined by the endpoint pairing: two pairs
cross iff their endpoints interleave.  In these forest cases there is
no extra isotopy datum; the regular neighbourhood of the forest is
unique once the cyclic order of the six boundary leaves is fixed.  The
unique all-crossing pairing is \(14|25|36\), so it is excluded.

Conversely, for any pairing whose interleaving graph is a forest, draw
one crossing for each interleaving pair.  The resulting
\(\Gamma(M)\) is a forest.  Face-boundary components are arcs between
boundary endpoints.  Push one such arc slightly into its adjacent face;
the pushed arc is disjoint from the diagram and separates the marked
points into two sets.  No matching interval crosses the pushed arc, so
each side contains an even number of marked points.  Hence the endpoints
of the face-boundary arc have opposite parity, exactly the inverse-label
condition for \(u\).  Thus admissible pairings are precisely all
pairings except \(14|25|36\).

Cells for \(u\):
\[
\begin{array}{lll}
v_0=12|34|56, & v_1=12|36|45, & v_2=14|23|56,\\
v_3=16|23|45, & v_4=16|25|34; &
\end{array}
\]
\[
\begin{array}{lll}
e_{01}=12|35|46, & e_{02}=13|24|56, & e_{13}=13|26|45,\\
e_{23}=15|23|46, & e_{04}=15|26|34, & e_{34}=16|24|35;
\end{array}
\]
\[
Q_{12}=13|25|46,\qquad Q_{14}=14|26|35,\qquad Q_{24}=15|24|36.
\]
Boundaries:
\[
\partial Q_{12}=e_{01}\cup e_{13}\cup e_{23}\cup e_{02},\quad
\partial Q_{14}=e_{01}\cup e_{13}\cup e_{34}\cup e_{04},\quad
\partial Q_{24}=e_{02}\cup e_{23}\cup e_{34}\cup e_{04}.
\]
The 1-skeleton is a theta graph from \(v_0\) to \(v_3\) with middle
vertices \(v_1,v_2,v_4\).  \(Q_{12}\) and \(Q_{14}\) glued along the
path through \(v_1\) form a disk whose boundary is the cycle through
\(v_2,v_4\); \(Q_{24}\) glues to that boundary.  Therefore the complex
is \(S^2\), not contractible.

\end{document}
